\section{结论与讨论}\label{ux7ed3ux8bbaux4e0eux8ba8ux8bba}

本文提出了一种面向多标签眼科疾病识别的可解释与隐私保护单域泛化框架，针对临床多疾病共存、多机构数据难以联合训练且跨机构部署存在域偏移、以及医疗数据隐私敏感等问题，构建了一种统一解决框架。核心贡献包括：（1）自适应标签-特征关联性构建模块（ARC），利用Transformer自/交叉注意力捕捉标签间依赖与特征-标签匹配关系，并通过一致性损失对齐双分支输出，有助于减少多标签识别中的误检和漏检；（2）双分支引导域不变特征提取模块（DFE），通过风格增强模拟潜在域偏移、CAM热图提供疾病区域先验引导，并结合空间/通道注意力与掩码过滤提取可迁移的疾病相关特征，在提升未知域泛化性能的同时增强模型可解释性；（3）自适应高斯差分隐私机制（AdaGDP），通过动态衰减噪声强度，在中等隐私预算（\(\epsilon \approx 10\)）下控制性能损失，实现隐私保护与模型可用性的平衡。系统性的实验验证了所提出方法的有效性。尽管取得一定进展，本框架仍存在局限：类别不平衡可能影响少数类疾病识别的稳定性，未来可结合重采样或损失加权进一步缓解；当前跨数据集泛化主要基于共有标签空间展开，未来可探索开放类别和更复杂分布外场景下的泛化能力。
